\section{Trimming the output for mass production}
\label{sec:trim}

The loose configuration of Section~\ref{sec:reco} stays the default for small diagnostic samples.
For the full-dataset production a second cut set, \fcmd{prod}, trims the output in four steps. It is
selected by the last argument of the macro (\fcmd{Fun4All_run3pp_KFP.C(...,"prod")}, i.e.\
\fcmd{macro/run_kfp.sh <nev> <dst> <outdir> 0 0 Ks,Lambda,D0 prod} or
\fcmd{condor/submit.sh <tag> <list> 0 prod}) and is always run without the crossing-study
channels. Each step was tested on the same 31 segments of run 79509 (\trimNev{} events).

\subsection{Trimming steps}
\begin{enumerate}\setlength\itemsep{2pt}
\item \textbf{No crossing-study output.} The \fcmd{*_noXmatch} files (\trimNoXMB\,MB, \trimNoXKB\,kB per
  event in the test) are not produced.
\item \textbf{Daughter impact parameter} (KFParticle \fcmd{setMinimumTrackPV_DCA}, 3D DCA, minimum over
  the PVs of the track's crossing, Section~\ref{sec:trackcuts}): $>0.05$\,cm for \Ks{} and \Lam{}
  daughters, $>0.003$\,cm (30\,$\mu$m) for \Dzero{} daughters; in this step also DIRA$>0.8$ for \Dzero.
\item \textbf{Slim format} (\fpath{tools/slim_kfp.C}, run inside the job after \fcmd{se->End()}): the
  64 trigger branches are not booked (\fcmd{getTriggerInfo(false)}, no \fcmd{TriggerRunInfoReco}; the
  trigger vectors remain per event in \fcmd{EventInfo}); redundant momentum branches are dropped
  (mother: $p_x,p_y,p_z,E,p,\sigma_p,\theta,y$, charge; daughters: $p_x,p_y,p_z,E,p,\sigma_p,\theta,y$,
  mass, $j_T$, $x,y,z$; and \fcmd{secondary_vertex_mass_pionPID}); the tree is rewritten with ZSTD
  compression and a 32\,kB basket size. The latter fixes a storage inefficiency of KFParticle, which
  books the covariance branches with a buffer of 6 or 21 \emph{bytes}, i.e.\ one basket per entry.
  For \Ks{} the four covariance branches took 640 bytes per candidate before and 260 after the fix
  (PV 121$\to$15, mother 171$\to$81, daughters 348$\to$163; \fpath{logs/branch_sizes_Ks.txt}).
\item \textbf{No covariance matrices} (mother 21, daughters $2\times21$, PV 6 elements).
\item \textbf{DIRA$>0.9$} for all three channels.
\end{enumerate}
The dropped kinematic quantities follow from the kept ones: $p=p_T\cosh\eta$, and the daughter mass
from its PDG code (needed e.g.\ for the swapped \Dzero{} mass). Duplicates across events are removed
with the key (\fcmd{Collision_BCO}, mass, \fcmd{track_1_pT}), because \fcmd{track_1_px} no longer exists.

\subsection{Size and signal}
\begin{table}[H]
\centering\footnotesize
\setlength\tabcolsep{4pt}
\caption{Cumulative effect of the trimming steps on the same \trimNev{} events
(\fpath{scripts/trim_study.C}). Candidates and sizes of the \Ks, \Lam, \Dzero{} trees; MB and kB per
event are file sizes of the merged ntuples (the \fcmd{*_noXmatch} files and the \fcmd{EventInfo}
file, \trimEvInfoKB\,kB/event in all versions, are not included). Signal: \Ks{} and \Lam+\aLam{}
yield with the offline QA cuts of Section~\ref{sec:qacuts} (Gaussian+pol2 fit, duplicates kept,
statistical errors).}
\label{tab:trim}
\begin{tabular}{rp{0.26\textwidth}rrrrrrrr}
\toprule
 & version & br. & \Ks{} cand. & \Lam{} cand. & \Dzero{} cand. & MB & kB/ev & \Ks{} signal & \Lam{} signal \\
\midrule
\input{tables/tab_trim.tex}
\bottomrule
\end{tabular}
\end{table}

\begin{itemize}\setlength\itemsep{1pt}
\item The three candidate trees shrink from \trimKBLoose{} to \trimKBDira\,kB per event (\trimFracDira\%
  of the loose size); with the crossing-study files of the diagnostic run included the reduction is a
  factor $\approx$\trimTotalReduction.
\item The largest single gains come from the daughter impact-parameter cut (version 0$\to$1:
  \trimFracProd\% of the loose size) and from the covariance branches (versions 1$\to$3: basket fix and
  removal, \trimKBProd$\to$\trimKBNoCov\,kB/event). The trigger block was only 0.6\% of the size.
\item With the QA cuts the \Ks{} and \Lam{} yields are unchanged within less than 1\% through all steps.
  Since the QA cuts already contain DIRA$>0.995$ and $|{\rm DCA}_{xy}|>0.02$\,cm, this shows that the
  trimming does not cut into that selection; it does not measure the loss for a looser final selection.
\item Signal removed by DIRA$>0.9$: with the other QA cuts but no DIRA cut, sideband-subtracted counts in
  the version before the change give \Ks{} $\bandLoKs\pm\ebandLoKs$ in $0.8<{\rm DIRA}\leq0.9$ against
  $\bandHiKs\pm\ebandHiKs$ above 0.9, and \Lam{} $\bandLoLambda\pm\ebandLoLambda$ against
  $\bandHiLambda\pm\ebandHiLambda$, i.e.\ no significant loss.
\end{itemize}

\caveat{The sizes in Table~\ref{tab:trim} are of merged files written by \fcmd{hadd} (zlib level 1).
The per-segment job outputs, which are what a production stores, measure 3.72, 1.76 and 0.62\,kB/event
for versions 1, 2 and 4 (\fpath{logs/trim_persegment.txt}), within 2--7\% of the merged values. The signal
loss of the impact-parameter cuts was not measured without the QA $|{\rm DCA}_{xy}|>0.02$\,cm cut, and
the \Dzero{} loss cannot be measured in this sample (no peak).}

\paragraph{Production estimate.} With \trimKBDira+\trimEvInfoKB\,kB per event, the $\approx5\times10^8$
events in the DST area at the time of the test (\nDSTfiles{} segments) give $\approx$\trimTBDira\,TB,
compared with \trimTBLoose\,TB for the loose configuration without the crossing-study files. The three
KFParticle modules need 63+71+94\,ms per event in the final configuration
(\fpath{logs/trim_persegment.txt}).
